\chapter{Chapter 18: Design a Voice-First Real-Time Agent (AI Drive-Thru / Phone Agent)}

\section*{Context}
\textbf{Problem statement:} Design a real-time voice agent that conducts natural conversations over phone or drive-thru intercom with sub-500ms voice-to-voice latency, handles interruptions gracefully, and maintains task state across multi-turn dialogues---processing 50K concurrent sessions with human-like responsiveness.

\subsection*{Real-World Examples}
\begin{itemize}
  \item A fast-food chain deploys AI drive-thru agents across 5,000 locations. Each agent handles 200 orders/day. The critical metric: voice-to-voice latency (time from when the customer stops speaking to when the agent starts responding). At 800ms, customers perceive the agent as ``slow and robotic'' and 30\% abandon mid-order. At 400ms, satisfaction scores match human operators. The system must handle interruptions---a customer changing their order mid-sentence (``Actually, make that a large---no, medium'') requires the agent to stop speaking, reparse, and respond within 500ms.
  \item A healthcare appointment scheduling system handles 50K concurrent calls during open enrollment. Callers frequently interrupt the agent (``I already said Tuesday'') and speak over confirmations. A naive pipeline (wait for silence $\rightarrow$ STT $\rightarrow$ LLM $\rightarrow$ TTS) produces 1.2s latency that causes callers to repeat themselves, creating a feedback loop of confusion. Switching to a streaming pipeline (STT starts transcribing during speech, LLM begins generating before STT finalizes) cuts latency to 450ms and reduces repeat-utterance rate by 60\%.
\end{itemize}

\bigskip\hrule\bigskip

\section*{Step 1: Understand the Problem and Establish Design Scope}

\subsection*{Simulated Dialogue}
\textbf{Interviewer:} Design a voice-first real-time agent for drive-thru ordering or phone-based customer service.\\
\textbf{Candidate:} Voice agents have a fundamentally different latency contract than text-based systems. In text chat, 2 seconds is fine. In voice, anything over 500ms feels unnatural---humans expect conversational pacing. I need to understand the voice pipeline architecture (STT$\rightarrow$LLM$\rightarrow$TTS vs.\ native multimodal), interruption handling requirements, and concurrent session targets.

\textbf{Candidate:} What's the scale and concurrency?\\
\textbf{Interviewer:} 50K concurrent voice sessions at peak. Each session lasts 2--5 minutes with 5--15 conversational turns.

\textbf{Candidate:} What latency targets?\\
\textbf{Interviewer:} Voice-to-voice latency (end of user speech to start of agent speech) must be $<$500ms at p95. Total response must begin within 500ms; full responses stream over 2--5 seconds.

\textbf{Candidate:} What quality metrics?\\
\textbf{Interviewer:} STT accuracy $>$95\% WER (word error rate $<$5\%). TTS naturalness MOS $>$4.0 (Mean Opinion Score, 1--5 scale). Task completion rate $>$90\% (e.g., order placed correctly). Interruption handling: agent must stop within 200ms of user barge-in.

\textbf{Candidate:} What's the deployment environment?\\
\textbf{Interviewer:} Phone calls via SIP/PSTN for customer service. Drive-thru intercom via audio stream. Both require handling background noise (kitchen noise, traffic, multiple speakers).

\subsection*{Functional Requirements}
\begin{itemize}
  \item Real-time bidirectional voice conversation: agent speaks while listening for interruptions.
  \item STT with streaming partial transcripts (not wait-for-end).
  \item LLM response generation triggered on partial transcripts (speculative execution).
  \item TTS with streaming audio output (first audio chunk before full response is generated).
  \item Barge-in detection: if the user speaks while the agent is speaking, the agent stops and listens.
  \item Task-state management: order state, appointment details, or issue context maintained across turns.
  \item Multi-language support for common scenarios (English + Spanish at minimum).
  \item Graceful escalation to human agent when confidence drops below threshold.
\end{itemize}

\subsection*{Non-Functional Requirements}
\begin{itemize}
  \item \textbf{Voice-to-voice latency:} p95 $<$ 500ms (user silence detected $\rightarrow$ first agent audio byte).
  \item \textbf{Concurrent sessions:} 50K peak.
  \item \textbf{STT accuracy:} WER $<$ 5\% on domain-specific vocabulary.
  \item \textbf{TTS naturalness:} MOS $>$ 4.0.
  \item \textbf{Barge-in response:} agent audio stops within 200ms of user voice detection.
  \item \textbf{Availability:} 99.95\% (voice downtime = lost revenue).
  \item \textbf{Cost:} $<$ \$0.05 per minute of conversation.
\end{itemize}

\subsection*{Back-of-the-Envelope Estimation}

\subsubsection*{Session Profile}
\begin{itemize}
  \item Concurrent sessions: 50K
  \item Average session duration: 3 minutes
  \item Turns per session: 10 (5 user, 5 agent)
  \item Total sessions/day: assuming average session = 3 min, 50K concurrent, avg utilization 40\%: \texttt{50K * 0.4 * (86400 / 180) = 9.6M sessions/day}
  \item Realistically (business hours, 12hr window): \texttt{50K * (12 * 3600 / 180) = 12M sessions/day}
\end{itemize}

\subsubsection*{Audio Data Volume}
\begin{itemize}
  \item Audio bitrate: 16kHz, 16-bit mono = 256 kbps = 32 KB/s
  \item Per session (3 min): \texttt{32 * 180 = 5.76 MB}
  \item Daily audio: \texttt{12M * 5.76 MB = 69 TB} (if stored; typically not stored long-term)
  \item Audio in transit: 50K sessions $\times$ 32 KB/s = 1.6 GB/s sustained bandwidth
\end{itemize}

\subsubsection*{STT Compute}
\begin{itemize}
  \item Whisper large-v3: processes 30s of audio in $\sim$1s on A100 GPU (30x real-time)
  \item For streaming STT (Deepgram, AssemblyAI): cloud API, no GPU management needed
  \item Self-hosted streaming STT: each GPU handles $\sim$30 concurrent streams
  \item GPUs for 50K concurrent: \texttt{50K / 30 = \textasciitilde{}1,667 GPUs} for STT alone (prohibitive self-hosted)
  \item Reality: use managed STT API (Deepgram: \$0.0043/min, AssemblyAI: \$0.0065/min)
  \item STT cost: \texttt{12M sessions * 3 min * \$0.0043 = \$154,800/day} $\leftarrow$ significant
\end{itemize}

\subsubsection*{LLM Compute}
\begin{itemize}
  \item Tokens per turn: $\sim$100 input (transcript + state), $\sim$50 output
  \item Turns per session: 5 agent turns
  \item Tokens per session: \texttt{5 * 150 = 750 tokens}
  \item Total tokens/day: \texttt{12M * 750 = 9B tokens}
  \item Using a fast 7B model (self-hosted, \$0.10/1M tokens): \texttt{9B * \$0.10/1M = \$900/day}
  \item Using GPT-4o-mini (\$0.15/1M input, \$0.60/1M output): $\sim$\$2,700/day
  \item LLM is NOT the cost bottleneck. STT and TTS are.
\end{itemize}

\subsubsection*{TTS Compute}
\begin{itemize}
  \item Agent speaks $\sim$15 seconds per turn, 5 turns = 75 seconds per session
  \item ElevenLabs: \$0.18/1K chars, $\sim$15 chars/second = $\sim$\$0.04/min
  \item Self-hosted (VITS/Bark): 1 GPU handles $\sim$100 concurrent streams at real-time factor 50x
  \item GPUs for 50K concurrent TTS: \texttt{50K / 100 = 500 GPUs} self-hosted
  \item Managed TTS cost: \texttt{12M * 1.25 min * \$0.04 = \$600K/day} $\leftarrow$ prohibitive at scale
  \item Self-hosted TTS is mandatory at this scale
\end{itemize}

\subsubsection*{Latency Budget (500ms voice-to-voice)}

\begin{tabular}{|p{0.35\linewidth}|p{0.08\linewidth}|p{0.47\linewidth}|}
\hline
\textbf{Stage} & \textbf{Budget} & \textbf{Notes} \\
\hline
Voice Activity Detection (VAD) & 50ms & Detect end-of-speech (silence threshold) \\
\hline
STT finalization (last partial $\rightarrow$ final) & 100ms & Streaming STT already has partial; final confirmation \\
\hline
Network (STT $\rightarrow$ LLM) & 20ms & Co-located services \\
\hline
LLM prefill + first token & 150ms & 7B model, $\sim$100 input tokens \\
\hline
LLM first token $\rightarrow$ TTS first chunk & 30ms & Streaming handoff \\
\hline
TTS first audio chunk generation & 100ms & First 200ms of audio generated \\
\hline
Network (TTS $\rightarrow$ user) & 50ms & Audio packet delivery \\
\hline
\textbf{Total} & \textbf{500ms} & Zero slack---every stage must hit budget \\
\hline
\end{tabular}

\subsubsection*{Cost Summary}

\begin{tabular}{|p{0.14\linewidth}|p{0.2\linewidth}|p{0.22\linewidth}|p{0.3\linewidth}|}
\hline
\textbf{Component} & \textbf{Self-Hosted (per day)} & \textbf{Managed API (per day)} & \textbf{Notes} \\
\hline
STT & \$12K (400 GPUs) & \$155K (Deepgram) & Self-host at scale \\
\hline
LLM & \$900 (30 GPUs) & \$2.7K (GPT-4o-mini) & Either works \\
\hline
TTS & \$15K (500 GPUs) & \$600K (ElevenLabs) & Must self-host \\
\hline
Infrastructure & \$5K (networking, storage) & \$2K & \\
\hline
\textbf{Total} & \textbf{$\sim$\$33K/day} & \textbf{$\sim$\$760K/day} & Self-hosted is 23x cheaper \\
\hline
\textbf{Per minute} & \textbf{\$0.0019} & \textbf{\$0.044} & \\
\hline
\end{tabular}

\bigskip\hrule\bigskip

\section*{Step 2: Propose High-Level Design and Get Buy-In}

\subsection*{Core Architecture Components}
\begin{itemize}
  \item \textbf{Voice Gateway:} Handles SIP/PSTN connections for phone calls, WebSocket/WebRTC for web-based voice, and audio stream connections for drive-thru intercoms. Manages audio codec negotiation, echo cancellation, and noise suppression.
  \item \textbf{VAD (Voice Activity Detection):} Continuously monitors the audio stream. Detects speech onset (user starts talking) and speech offset (user stops talking). Also detects barge-in (user talks while agent is speaking). Low-latency requirement: $<$50ms detection.
  \item \textbf{Streaming STT:} Converts audio to text in real-time with partial (interim) transcripts. Partial transcripts arrive every 100--300ms as the user speaks. Final transcript arrives 100--200ms after speech ends.
  \item \textbf{Dialogue Manager:} Maintains conversation state (order items, customer info, issue context). Receives transcripts, determines next action (respond, ask clarification, execute tool call), manages multi-turn context window.
  \item \textbf{LLM Engine:} Generates agent responses based on transcript + dialogue state + system prompt. Must support streaming output (token-by-token). Optimized for low TTFT ($<$150ms).
  \item \textbf{Streaming TTS:} Converts LLM text output to audio in real-time. Starts generating audio from the first sentence fragment, doesn't wait for full response. Streams audio chunks back to the user.
  \item \textbf{Barge-In Controller:} Coordinates interruption handling. When VAD detects user speech during agent TTS playback: (1) immediately stop TTS audio, (2) cancel in-flight LLM generation, (3) route new audio to STT, (4) update dialogue state with interruption context.
  \item \textbf{Session State Store:} Low-latency key-value store (Redis) holding per-session state: conversation history, order state, customer metadata, current turn status.
  \item \textbf{Audio Transport Layer:} WebRTC for low-latency bidirectional audio. Jitter buffer management, packet loss concealment, adaptive bitrate.
  \item \textbf{Observability:} Voice-to-voice latency histograms, STT WER sampling, TTS MOS sampling, barge-in rate, task completion rate, escalation rate.
\end{itemize}

\subsection*{Pipeline A: Session Establishment Pipeline}
\begin{enumerate}
  \item Incoming call (SIP) or audio connection (WebSocket) arrives at \textbf{Voice Gateway}.
  \item Gateway negotiates codec (Opus for WebRTC, G.711 for PSTN), establishes bidirectional audio stream.
  \item New session created in \textbf{Session State Store} with initial context (caller ID, time, location for drive-thru).
  \item \textbf{VAD} begins monitoring inbound audio stream.
  \item \textbf{TTS} generates and plays the greeting: ``Welcome to [restaurant], what can I get for you?''
  \item System enters listening state.
\end{enumerate}

\subsection*{Pipeline B: Turn-by-Turn Conversation Pipeline}
\begin{enumerate}
  \item \textbf{User speaks.} VAD detects speech onset $\rightarrow$ \textbf{Streaming STT} begins transcribing.
  \item STT emits partial transcripts every 200ms: ``I'd like a\ldots'' $\rightarrow$ ``I'd like a large\ldots'' $\rightarrow$ ``I'd like a large combo number three.''
  \item \textbf{Speculative execution (optional):} after the first 2--3 partial transcripts, the \textbf{Dialogue Manager} begins pre-processing likely intents. If ``I'd like a large'' strongly predicts an order, pre-load the menu context.
  \item VAD detects end-of-speech (300ms silence) $\rightarrow$ STT emits final transcript.
  \item \textbf{Dialogue Manager} receives final transcript, resolves it against conversation state:
  \begin{itemize}
    \item Extracts intent (order item, modify order, ask question, confirm).
    \item Executes any tool calls (add item to order, check availability).
    \item Constructs LLM prompt: system prompt + menu/policy context + conversation history + current transcript.
  \end{itemize}
  \item \textbf{LLM Engine} generates response with streaming output. First token in $<$150ms.
  \item \textbf{Streaming TTS} receives first sentence fragment from LLM, begins audio synthesis. First audio chunk ready in $<$100ms.
  \item Audio streams to user. Agent is ``speaking.''
  \item \textbf{If user interrupts (barge-in):} VAD detects user speech $\rightarrow$ \textbf{Barge-In Controller} triggers:
  \begin{itemize}
    \item TTS audio playback stops within 200ms.
    \item LLM generation cancelled.
    \item STT begins transcribing the interruption.
    \item Dialogue Manager receives the interruption transcript with context: ``interrupted agent mid-response.''
    \item New LLM call generated with interruption context.
  \end{itemize}
  \item Turn completes. State updated. Next turn begins.
\end{enumerate}

\subsection*{Mermaid Architecture Diagram}

% Mermaid diagram (not natively supported in LaTeX)
\begin{verbatim}
flowchart LR
  caller[Caller Phone/Intercom] <-->|audio| voiceGW[VoiceGateway SIP/WebRTC]
  voiceGW --> vad[VAD SileroVAD]
  voiceGW <-->|audio stream| transport[AudioTransport Jitter Buffer]

  vad -->|speech onset| stt[StreamingSTT Whisper/Deepgram]
  vad -->|end of speech| stt
  vad -->|barge-in| bargeIn[BargeInController]

  stt -->|partial transcripts| dialogMgr[DialogueManager]
  stt -->|final transcript| dialogMgr

  dialogMgr --> sessionStore[SessionStateStore Redis]
  dialogMgr --> toolExec[ToolExecutor Order/Lookup]
  dialogMgr --> llm[LLMEngine 7B Streaming]

  llm -->|streaming tokens| tts[StreamingTTS VITS/XTTS]
  tts -->|audio chunks| transport
  transport --> caller

  bargeIn -->|stop| tts
  bargeIn -->|cancel| llm
  bargeIn -->|interrupt context| dialogMgr

  observability[Observability Latency WER MOS BargeIn] --> stt
  observability --> llm
  observability --> tts
  observability --> vad
\end{verbatim}

\bigskip\hrule\bigskip

\section*{Step 3: Design Deep Dive}

Focus areas:
\begin{enumerate}
  \item \textbf{End-to-End Voice Latency Pipeline}
  \item \textbf{Interruption Handling and Turn-Taking}
  \item \textbf{Pipeline vs.\ Native Multimodal Architecture}
\end{enumerate}

\subsection*{1) End-to-End Voice Latency Pipeline}

The 500ms budget is the single most important constraint. Every millisecond wasted in one stage must be recovered in another. The pipeline must be designed for streaming at every stage---no stage waits for the previous stage to complete.

\textbf{The streaming pipeline principle:}

In a naive (batch) pipeline: user finishes speaking $\rightarrow$ STT processes full utterance $\rightarrow$ LLM receives full text $\rightarrow$ generates full response $\rightarrow$ TTS synthesizes full audio $\rightarrow$ playback begins. Total: 2--3 seconds.

In a streaming pipeline: STT emits partial transcripts while user speaks $\rightarrow$ LLM begins generating from partial context $\rightarrow$ TTS synthesizes audio from first sentence fragment $\rightarrow$ first audio byte plays before LLM finishes generating. Total: 400--500ms from end-of-speech to first audio.

\textbf{Per-stage latency optimization:}

\textbf{Stage 1: VAD (Voice Activity Detection) --- 50ms budget}

\begin{tabular}{|p{0.15\linewidth}|p{0.08\linewidth}|p{0.08\linewidth}|p{0.1\linewidth}|p{0.47\linewidth}|}
\hline
\textbf{VAD Model} & \textbf{Latency} & \textbf{Accuracy} & \textbf{Resource} & \textbf{Notes} \\
\hline
Silero VAD & 10--30ms & 95\%+ & CPU only & \textbf{Production default.} Lightweight, runs on any CPU. \\
\hline
WebRTC VAD (Google) & 5--10ms & 90\% & CPU only & Very fast but higher false-positive rate in noisy environments \\
\hline
Energy-based (RMS threshold) & 1--5ms & 80\% & CPU only & Too many false triggers in drive-thru noise \\
\hline
Custom CNN VAD & 20--50ms & 97\% & GPU & Overkill unless noise environment is extreme \\
\hline
\end{tabular}

\bigskip
End-of-speech detection: VAD must decide ``the user has stopped talking'' quickly but not prematurely. Use a silence duration threshold:
\begin{itemize}
  \item 300ms silence $\rightarrow$ end-of-speech (default for transactional conversations)
  \item 500ms silence $\rightarrow$ end-of-speech (for more deliberate speakers)
  \item Adaptive: start at 300ms, increase to 500ms if the user has been pausing mid-sentence (detected via partial transcript analysis)
\end{itemize}

\textbf{Stage 2: Streaming STT --- 100ms budget (from end-of-speech to final transcript)}

The STT has been running throughout the user's speech, emitting partial transcripts. After VAD signals end-of-speech, the STT needs only $\sim$100ms to finalize the transcript (commit the last few words, apply language model correction).

{\small
\begin{tabular}{|p{0.14\linewidth}|p{0.07\linewidth}|p{0.08\linewidth}|p{0.1\linewidth}|p{0.09\linewidth}|p{0.1\linewidth}|p{0.15\linewidth}|}
\hline
\textbf{STT System} & \textbf{Streaming} & \textbf{WER (clean)} & \textbf{WER (noisy)} & \textbf{Latency (partial)} & \textbf{Latency (final)} & \textbf{Cost} \\
\hline
Deepgram Nova-2 & Yes & 3.2\% & 8.5\% & 200ms & 100ms after EOS & \$0.0043/min \\
\hline
Google Cloud STT v2 & Yes & 3.8\% & 9.1\% & 250ms & 150ms after EOS & \$0.006/min \\
\hline
Whisper large-v3 (streaming via whisper.cpp) & Semi (chunked) & 2.9\% & 7.8\% & 500ms chunks & 300ms after EOS & Self-hosted: $\sim$\$0.001/min \\
\hline
AssemblyAI Universal & Yes & 3.5\% & 8.0\% & 200ms & 100ms after EOS & \$0.0065/min \\
\hline
Fine-tuned Whisper (domain vocab) & Semi & 2.1\% & 5.2\% & 500ms chunks & 300ms after EOS & Self-hosted: $\sim$\$0.001/min \\
\hline
\end{tabular}
}

\bigskip
For $<$500ms voice-to-voice: managed streaming STT (Deepgram) provides the best latency profile. Self-hosted Whisper requires chunked streaming (500ms chunks), which adds 200ms to the pipeline.

\textbf{Domain vocabulary adaptation:}

Drive-thru menus have domain-specific terms (``McFlurry,'' ``Whopper,'' ``venti''). STT systems misrecognize these. Fix:
\begin{itemize}
  \item \textbf{Hot-word boosting:} Deepgram/Google Cloud support keyword boosting---provide a list of menu item names with boosted recognition probability. +15\% accuracy on menu items, zero latency impact.
  \item \textbf{Fine-tuning:} fine-tune Whisper on 100 hours of drive-thru audio with transcripts. WER drops from 8.5\% to 5.2\% on domain vocab. Requires labeled training data.
  \item \textbf{Post-correction:} after STT, run a fast string-matching pass to correct common misrecognitions: ``large MacFlurry'' $\rightarrow$ ``large McFlurry.'' Adds $<$5ms.
\end{itemize}

\textbf{Stage 3: LLM Generation --- 150ms budget (to first token)}

The LLM must produce the first output token within 150ms of receiving the final transcript. This constrains model size and optimization stack.

{\small
\begin{tabular}{|p{0.14\linewidth}|p{0.05\linewidth}|p{0.1\linewidth}|p{0.06\linewidth}|p{0.12\linewidth}|p{0.1\linewidth}|p{0.2\linewidth}|}
\hline
\textbf{Model} & \textbf{Size} & \textbf{TTFT (100 tok in)} & \textbf{TPOT} & \textbf{Quality (task compl.)} & \textbf{Self-hosted GPU} & \textbf{Notes} \\
\hline
GPT-4o-mini (API) & Unk. & 200--400ms & 15ms & 94\% & N/A & API latency adds 50--100ms network \\
\hline
Llama 3 8B (INT4, vLLM) & 8B & 80ms & 8ms & 88\% & $1\times$ A100 & \textbf{Production default} \\
\hline
Llama 3 8B + speculative & 8B+1B & 80ms & 3ms & 88\% & $1\times$ A100 & Fastest decode \\
\hline
Phi-3 3.8B (INT4) & 3.8B & 40ms & 5ms & 82\% & $1\times$ A10 & Budget option \\
\hline
Mistral 7B (INT4) & 7B & 75ms & 7ms & 86\% & $1\times$ A100 & Good alternative \\
\hline
\end{tabular}
}

\bigskip
For voice agents, a self-hosted 8B model with speculative decoding is optimal: 80ms TTFT, 3ms TPOT, fitting the 150ms budget with headroom. API-based models add 50--100ms network latency that's hard to recover.

\textbf{Prompt design for voice agents:}

Voice agent prompts differ from text agent prompts:
\begin{itemize}
  \item Responses must be SHORT. A 200-word text response is fine; a 200-word spoken response takes 60 seconds and the caller hangs up.
  \item Responses must be SPEAKABLE. No bullet points, URLs, or code. No parenthetical remarks.
  \item System prompt includes: ``Keep responses to 1--2 sentences. Use simple, conversational language. Never list more than 3 items. If the order is ambiguous, ask one clarifying question.''
  \item Token budget: cap output at 50 tokens ($\sim$15 seconds of speech). Longer responses get truncated.
\end{itemize}

\textbf{Stage 4: Streaming TTS --- 100ms budget (to first audio chunk)}

TTS begins synthesizing as soon as the first complete sentence or clause arrives from the LLM. It doesn't wait for the full response.

{\small
\begin{tabular}{|p{0.13\linewidth}|p{0.07\linewidth}|p{0.05\linewidth}|p{0.09\linewidth}|p{0.08\linewidth}|p{0.12\linewidth}|p{0.24\linewidth}|}
\hline
\textbf{TTS System} & \textbf{Streaming} & \textbf{MOS} & \textbf{Latency (1st chunk)} & \textbf{Languages} & \textbf{Cost} & \textbf{Notes} \\
\hline
ElevenLabs Turbo v2 (API) & Yes & 4.5 & 200--400ms & 30+ & \$0.18/1K chars & High quality but API latency \\
\hline
XTTS v2 (self-hosted) & Yes & 4.2 & 80ms & 17 & Self-hosted: $\sim$\$0.001/min & \textbf{Production default} \\
\hline
VITS (self-hosted) & Yes & 3.8 & 40ms & Per-language & Self-hosted: $\sim$\$0.0005/min & Fastest, lower quality \\
\hline
Bark (self-hosted) & No (batch) & 4.3 & 2,000ms & 10+ & Self-hosted & Too slow for real-time \\
\hline
Azure Neural TTS (API) & Yes & 4.4 & 150--300ms & 100+ & \$16/1M chars & Good quality, API latency \\
\hline
\end{tabular}
}

\bigskip
For $<$500ms pipeline: self-hosted XTTS v2 with 80ms first-chunk latency. Trade-off: slightly lower naturalness (MOS 4.2 vs 4.5) but meets the latency budget.

\textbf{Sentence-level streaming:}

The LLM generates tokens one at a time. TTS needs at least a sentence fragment to produce natural-sounding audio. Strategy:
\begin{itemize}
  \item Buffer LLM tokens until a sentence boundary (period, question mark) or a pause-worthy clause boundary (comma followed by a conjunction).
  \item Minimum fragment: 5 words. Shorter fragments produce choppy audio.
  \item Maximum wait before forcing synthesis: 500ms of buffered tokens. If the LLM is slow, synthesize whatever is available.
\end{itemize}

\subsubsection*{End-to-End Latency Trade-off Table}

{\small
\begin{tabular}{|p{0.15\linewidth}|p{0.1\linewidth}|p{0.13\linewidth}|p{0.1\linewidth}|p{0.12\linewidth}|p{0.22\linewidth}|}
\hline
\textbf{Pipeline Mode} & \textbf{V2V Latency} & \textbf{Quality} & \textbf{Complexity} & \textbf{Cost} & \textbf{Best For} \\
\hline
Batch (wait for each stage) & 2,000--3,000ms & Highest (no partial errors) & Low & Low & Voicemail transcription, non-interactive \\
\hline
Streaming STT + batch LLM + batch TTS & 1,200--1,800ms & High & Medium & Medium & Low-priority phone IVR \\
\hline
Full streaming (STT$\rightarrow$LLM$\rightarrow$TTS all streaming) & 400--500ms & Good (occasional partial-transcript errors) & High & Medium & \textbf{Production default for real-time agents} \\
\hline
Speculative pipeline (start LLM on partial transcript) & 300--400ms & Lower (may generate wrong response) & Very high & Higher (wasted compute) & Ultra-low-latency premium \\
\hline
\end{tabular}
}

\subsection*{2) Interruption Handling and Turn-Taking}

Humans interrupt each other constantly in conversation. A voice agent that can't handle interruptions feels robotic and frustrating. Barge-in handling is the second most critical feature after latency.

\textbf{Barge-in detection pipeline:}

\begin{enumerate}
  \item While TTS is playing (agent speaking), VAD continues monitoring the inbound audio stream.
  \item \textbf{Challenge:} the agent's own audio (played through a speaker/earpiece) feeds back into the microphone, creating an echo. VAD must distinguish the user's voice from the echo of the agent's voice.
  \item \textbf{Acoustic Echo Cancellation (AEC):} subtract the known agent audio signal from the incoming audio before VAD analysis. This is standard in telephony hardware (PSTN handles it) and WebRTC (built-in AEC). Drive-thru intercoms require custom AEC tuning for the acoustic environment.
  \item After AEC, if VAD detects speech energy above threshold $\rightarrow$ \textbf{barge-in detected.}
\end{enumerate}

\textbf{Barge-in response sequence (must complete in $<$200ms):}

\begin{tabular}{|p{0.04\linewidth}|p{0.4\linewidth}|p{0.08\linewidth}|p{0.36\linewidth}|}
\hline
\textbf{Step} & \textbf{Action} & \textbf{Latency} & \textbf{Notes} \\
\hline
1 & VAD detects speech during agent playback & 30ms & After AEC \\
\hline
2 & Send stop signal to TTS audio output & 10ms & Immediate audio mute \\
\hline
3 & Cancel in-flight LLM generation & 10ms & Abort streaming response \\
\hline
4 & Record how much of the agent response was played & 5ms & For context: ``agent said X before interrupted'' \\
\hline
5 & Route inbound audio to STT & 10ms & STT begins transcribing interruption \\
\hline
6 & Update dialogue state with interruption & 5ms & Flag: partial response delivered \\
\hline
\textbf{Total} & & \textbf{70ms} & Well within 200ms budget \\
\hline
\end{tabular}

\bigskip
\textbf{Turn-taking state machine:}

\begin{verbatim}
States: LISTENING -> PROCESSING -> SPEAKING -> INTERRUPTED

LISTENING: VAD monitors. On speech_onset -> route to STT.
           On end_of_speech -> transition to PROCESSING.
PROCESSING: STT finalizes. LLM generates.
            On first_audio_chunk -> transition to SPEAKING.
SPEAKING: TTS plays. On barge_in -> transition to INTERRUPTED.
          On playback_complete -> transition to LISTENING.
INTERRUPTED: Stop TTS. Cancel LLM. STT transcribes interruption.
             On end_of_speech -> transition to PROCESSING
             (with interruption context).
\end{verbatim}

\textbf{Context handling after interruption:}

When the user interrupts, the dialogue manager must know:
\begin{itemize}
  \item What the agent was saying when interrupted (for context, the user heard part of it).
  \item What the user said during/after the interruption.
  \item Whether the interruption is a correction (``no, I said medium''), an addition (``and also a drink''), or a redirect (``actually, cancel that'').
\end{itemize}

Prompt augmentation after interruption:
\begin{verbatim}
[Previous agent response, spoken up to: "Your total comes to $12.50.
  Would you like to add--"]
[User interrupted with: "Actually, remove the fries."]
[Instruction: The user interrupted your response. Acknowledge the
  interruption, process their request, and continue.]
\end{verbatim}

\textbf{False barge-in handling:}

Not all detected speech during agent playback is an intentional interruption:
\begin{itemize}
  \item Background noise (other people talking, traffic)
  \item The user saying ``uh-huh'' or ``yeah'' as backchannel acknowledgment
  \item Someone else at the drive-thru window talking
\end{itemize}

Strategy:
\begin{itemize}
  \item \textbf{Minimum speech duration:} don't trigger barge-in for speech bursts $<$ 300ms. Short bursts are likely backchannel or noise.
  \item \textbf{Confidence threshold:} if STT produces only filler words (``uh,'' ``hmm,'' ``yeah'') with high confidence, treat as backchannel, not interruption. Continue agent playback.
  \item \textbf{Two-stage detection:} on first speech detection, lower TTS volume by 50\% for 500ms. If speech continues past 500ms, fully stop TTS and process as interruption. If speech stops, restore TTS volume. This creates a graceful ``pause'' rather than a jarring stop.
\end{itemize}

\subsubsection*{Interruption Handling Trade-off Table}

\begin{tabular}{|p{0.2\linewidth}|p{0.12\linewidth}|p{0.12\linewidth}|p{0.2\linewidth}|p{0.12\linewidth}|}
\hline
\textbf{Strategy} & \textbf{False Barge-In Rate} & \textbf{Missed Barge-In Rate} & \textbf{User Experience} & \textbf{Complexity} \\
\hline
Immediate stop on any speech & 25\% (noise triggers) & 0\% & Choppy---agent stops for every noise & Low \\
\hline
300ms minimum speech + backchannel filter & 5\% & 3\% (misses very short corrections) & Good for most environments & Medium \\
\hline
Two-stage (volume reduction $\rightarrow$ full stop) & 3\% & 2\% & \textbf{Most natural} & High \\
\hline
No barge-in (wait for agent to finish) & 0\% & 100\% & Terrible---caller must wait & Trivial \\
\hline
\end{tabular}

\subsection*{3) Pipeline vs.\ Native Multimodal Architecture}

The traditional voice agent uses a cascade pipeline: STT $\rightarrow$ LLM $\rightarrow$ TTS. An emerging alternative: native speech-to-speech models that process audio directly without text intermediation.

\textbf{Pipeline architecture (STT $\rightarrow$ LLM $\rightarrow$ TTS):}

\begin{verbatim}
Audio In -> STT -> Text -> LLM -> Text -> TTS -> Audio Out
\end{verbatim}

Each stage is independently optimizable and debuggable. STT, LLM, and TTS can be swapped independently. Text intermediation enables logging, filtering, and safety checks on the actual words.

\textbf{Native multimodal architecture (Audio $\rightarrow$ Model $\rightarrow$ Audio):}

\begin{verbatim}
Audio In -> Multimodal Model -> Audio Out
\end{verbatim}

A single model (e.g., GPT-4o voice mode, Moshi, SpeechGPT) takes audio tokens as input and produces audio tokens as output. No text intermediation.

{\small
\begin{tabular}{|p{0.14\linewidth}|p{0.18\linewidth}|p{0.18\linewidth}|p{0.38\linewidth}|}
\hline
\textbf{Dimension} & \textbf{Pipeline (STT$\rightarrow$LLM$\rightarrow$TTS)} & \textbf{Native Multimodal} & \textbf{Notes} \\
\hline
\textbf{Voice-to-voice latency} & 400--500ms (sum of stages) & 200--300ms (single model) & Multimodal eliminates inter-stage latency \\
\hline
\textbf{STT accuracy dependency} & Yes---LLM only sees STT output; STT errors propagate & No---model hears raw audio, can understand despite noise & Major advantage for noisy environments \\
\hline
\textbf{LLM reasoning quality} & High---text-mode LLMs are well-optimized & Lower---speech-to-speech models are less capable at reasoning & Pipeline wins for complex tasks \\
\hline
\textbf{TTS naturalness} & Depends on TTS model (MOS 3.8--4.5) & Very high---model controls prosody, emotion, pacing natively (MOS 4.5+) & Multimodal produces more human-like speech \\
\hline
\textbf{Debuggability} & High---inspect text at each stage & Low---audio-in/audio-out is a black box & Critical for production debugging \\
\hline
\textbf{Safety/filtering} & Easy---text-based content filtering between LLM and TTS & Hard---must analyze audio output or add separate ASR for monitoring & Compliance requirement \\
\hline
\textbf{Language/voice control} & Separate TTS voice selection, style control & Limited---model has its own voice(s) & Pipeline offers more customization \\
\hline
\textbf{Cost} & Sum of 3 models & Single model (but larger, GPU-intensive) & Roughly comparable at scale \\
\hline
\textbf{Maturity} & Production-proven (2+ years) & Early production (2024--2025) & Pipeline is safer choice today \\
\hline
\textbf{Modularity} & Swap any component independently & Monolithic---upgrade is all-or-nothing & Pipeline wins for iteration speed \\
\hline
\textbf{Emotion/tone} & TTS must be explicitly instructed (SSML) & Model naturally modulates tone based on context & Multimodal feels more human \\
\hline
\end{tabular}
}

\bigskip
\textbf{Recommended: Pipeline with multimodal as a staged migration.}

Start with the pipeline architecture (STT$\rightarrow$LLM$\rightarrow$TTS) for production reliability, debuggability, and safety filtering. Monitor multimodal model maturity. Pilot multimodal on low-risk, high-volume scenarios (drive-thru greetings, simple order confirmations) while keeping pipeline for complex interactions (order modifications, complaints, escalations).

\textbf{Hybrid architecture (emerging best practice):}

\begin{verbatim}
Audio In -> STT -> Text -> LLM (text) -> Text -> TTS -> Audio Out
                    |
              Audio In -> Multimodal Model -> Audio Out
              (fallback for noisy audio where STT fails)
\end{verbatim}

Use the pipeline by default. If STT confidence drops below 70\% (noisy environment), fall back to the multimodal model that can handle audio directly. This gives the best of both worlds.

\subsubsection*{Architecture Strategy Trade-off Table}

{\small
\begin{tabular}{|p{0.15\linewidth}|p{0.08\linewidth}|p{0.1\linewidth}|p{0.1\linewidth}|p{0.1\linewidth}|p{0.1\linewidth}|p{0.2\linewidth}|}
\hline
\textbf{Architecture} & \textbf{Latency} & \textbf{Quality} & \textbf{Safety} & \textbf{Debug.} & \textbf{Maturity} & \textbf{Best For} \\
\hline
Pipeline (STT$\rightarrow$LLM$\rightarrow$TTS) & 400--500ms & Good & High (text filter) & High & Production & \textbf{Default for most deployments} \\
\hline
Native multimodal & 200--300ms & Higher nat. & Low (audio black box) & Low & Early & Greenfield projects with lower compliance \\
\hline
Hybrid (pipeline + multimodal fallback) & 400ms default, 250ms fallback & Best overall & Medium & Medium & & \textbf{Target architecture (2025+)} \\
\hline
Pipeline with speculative LLM & 300--400ms & Good & High & High & Exper. & Latency-critical premium tier \\
\hline
\end{tabular}
}

\subsection*{Production Bottlenecks and Mitigations}

\begin{itemize}
  \item \textbf{STT latency spike under load}
  \begin{itemize}
    \item Symptom: managed STT API p95 latency jumps from 100ms to 500ms during peak hours, pushing total voice-to-voice latency to 900ms.
    \item Mitigation: multi-provider STT failover. Primary: Deepgram. Failover: Google Cloud STT. If primary latency exceeds 200ms for 3 consecutive calls, route to failover. Additionally: self-host a Whisper instance pool for fallback (higher latency but no rate limits).
  \end{itemize}

  \item \textbf{Background noise false-triggers VAD continuously}
  \begin{itemize}
    \item Symptom: drive-thru kitchen noise triggers speech detection, causing constant barge-in on agent speech.
    \item Mitigation: deploy a noise-robust VAD trained on drive-thru audio (fine-tune Silero VAD on 1,000 hours of drive-thru ambient noise). Additionally: use directional microphone hardware at the order point to reduce ambient noise pickup.
  \end{itemize}

  \item \textbf{TTS audio glitches during streaming}
  \begin{itemize}
    \item Symptom: audio artifacts (clicks, pops, gaps) at sentence boundaries when TTS switches from one synthesized fragment to the next.
    \item Mitigation: overlap-add blending at fragment boundaries. Generate 50ms of overlapping audio between consecutive fragments and crossfade. This smooths transitions. Additionally: buffer two fragments ahead before starting playback (adds 100ms to initial latency but eliminates all audible gaps).
  \end{itemize}

  \item \textbf{Dialogue state corruption on rapid interruptions}
  \begin{itemize}
    \item Symptom: user interrupts twice in quick succession (``no---wait---actually, make it a large''), and the dialogue manager loses track of the current order state.
    \item Mitigation: atomic state updates with version numbers. Each state change is a versioned transaction. Interrupted LLM generations that haven't committed a state change are rolled back. Only fully processed turns update the state.
  \end{itemize}

  \item \textbf{Caller on hold / silence handling}
  \begin{itemize}
    \item Symptom: caller goes silent for 10+ seconds (putting phone down, talking to someone else). System times out and disconnects.
    \item Mitigation: implement a ``patience protocol.'' After 5s silence: ``Take your time, I'm still here.'' After 15s: ``Are you still there?'' After 30s: ``I'll keep your order saved. Call back when you're ready.'' Don't hard-disconnect until 60s.
  \end{itemize}
\end{itemize}

\subsection*{Advanced Techniques}
\begin{itemize}
  \item \textbf{Speculative response generation:} while the user is still speaking (mid-utterance), the system predicts the likely complete utterance from partial transcripts and begins generating a response speculatively. If the prediction is correct ($>$70\% of the time for common patterns like ``I'd like a [menu\_item]''), the response is already partially generated by the time the user finishes. Saves 100--200ms. Risk: wasted compute on wrong predictions ($\sim$30\% waste rate).
  \item \textbf{Emotion detection for escalation:} analyze audio features (pitch, speaking rate, volume) to detect caller frustration. If frustration score exceeds threshold, proactively offer: ``I want to make sure I get this right for you. Would you like me to connect you with a team member?'' Early escalation improves CSAT vs.\ letting the caller struggle with the AI.
  \item \textbf{Contextual pre-loading:} for drive-thru, pre-load the current menu, promotions, and item availability into the LLM context before the first customer interaction. For phone agents, use caller ID to pre-load account history. Reduces per-turn context assembly time to near-zero.
\end{itemize}

\bigskip\hrule\bigskip

\section*{Step 4: Wrap Up}

\subsection*{Final Design Summary}

The system achieves sub-500ms voice-to-voice latency through a fully streaming pipeline: Silero VAD for end-of-speech detection (50ms), streaming STT with domain-adapted vocabulary (100ms finalization), a self-hosted 8B LLM with speculative decoding (80ms TTFT), and self-hosted XTTS v2 for streaming TTS (80ms first chunk). Barge-in handling uses a two-stage detection strategy (volume reduction then full stop) that achieves 3\% false positive rate while catching 98\% of real interruptions. At 50K concurrent sessions, self-hosted STT and TTS at \$0.002/minute is 23x cheaper than managed APIs. The architecture is pipeline-based (STT$\rightarrow$LLM$\rightarrow$TTS) for production maturity and debuggability, with a planned migration path to hybrid pipeline + native multimodal for noisy environments.

\subsection*{Wrap-Around Concerns}
\begin{itemize}
  \item \textbf{Risk:} 500ms latency budget has zero slack. Any single stage regression (STT provider latency spike, GPU contention on LLM, TTS audio buffer underrun) breaks the SLO. Mitigation: per-stage latency monitoring with circuit breakers. If any stage exceeds 1.5x budget, degrade gracefully (shorter responses, skip speculative execution, use lower-quality but faster TTS).
  \item \textbf{Risk:} Background noise in drive-thru environments degrades STT accuracy below 95\%. Mitigation: domain-adapted STT model, hot-word boosting, and multimodal fallback for low-confidence transcripts.
  \item \textbf{Risk:} Barge-in handling with AEC fails in novel acoustic environments (new intercom hardware, outdoor seating areas). Mitigation: per-location AEC calibration during installation. Monitor barge-in false-positive rate per location; alert if $>$10\%.
  \item \textbf{Risk:} Self-hosted TTS and STT require 900+ GPUs at peak, creating significant infrastructure operational burden. Mitigation: auto-scaling with spot instances for burst. Maintain 60\% baseline on reserved instances, scale 40\% on spot.
  \item \textbf{Risk:} Dialogue state errors cause wrong orders (customer says ``large'' but system records ``small''). Mitigation: explicit confirmation turns for high-value items. ``Got it---one large combo number three. Is that correct?'' Confirmation rate should be $>$95\%.
\end{itemize}

\bigskip\hrule\bigskip

\section*{Step 5: Interviewer Curveballs}

\subsection*{Interviewer Follow-ups \& Frequently Asked Questions (FAQ)}

\textbf{Q:} The customer says ``I'd like a large number three with no pickles and a Coke---oh wait, make that a Sprite.'' Your system must handle this single utterance. Walk through exactly what happens.

\textbf{The Ideal Answer:} This is a self-correction within a single turn---not a multi-turn dialogue issue. The STT transcribes the full utterance as one text. The LLM must parse the self-correction. Prompt design handles this: the system prompt instructs ``Customers may change their mind mid-sentence. The LAST stated preference overrides earlier ones.'' The LLM parses: (1) large \#3, (2) no pickles, (3) Coke, (4) correction: Sprite replaces Coke. Final order: large \#3, no pickles, Sprite. The dialogue manager commits this as a single state update. Confirmation: ``A large number three, no pickles, with a Sprite. Anything else?''

\textbf{Deep Dive \& Trade-offs:}
\begin{itemize}
  \item The LLM must be robust to self-corrections. Include few-shot examples of self-corrections in the system prompt.
  \item If STT segments this as two utterances (pause between ``Coke'' and ``oh wait''), the system processes two turns. The second turn (``make that a Sprite'') is a modification. Both paths work but the two-turn path adds 500ms.
  \item Edge case: ``no pickles'' is a negation. LLMs sometimes struggle with negations in structured extraction. Include explicit negation examples in few-shot.
\end{itemize}

\textbf{The Interviewer's Rubric:} Checking understanding of real-world speech patterns and how they map to dialogue management.

\bigskip\hrule\bigskip

\textbf{Q:} Your STT transcribes ``I'd like a four-piece nugget'' but the customer actually said ``forty-piece nugget.'' This is a real order---how do you prevent costly misrecognitions?

\textbf{The Ideal Answer:} Numerals in noisy environments are the \#1 STT failure mode. Defense layers: (1) \textbf{Hot-word boosting:} if the menu has ``4-piece'' and ``40-piece'' but no ``forty-piece'' as a string, STT may default to the more common ``four-piece.'' Boost all menu item names including their numeric variants. (2) \textbf{Menu-aware post-processing:} after STT, validate the transcript against the actual menu. ``Four-piece nugget'' is a valid menu item but so is ``forty-piece nugget.'' If ambiguous, use acoustic confidence scores on the numeral. (3) \textbf{Confirmation protocol:} for items where the quantity matters (and thus the price difference is large---\$3 vs \$12), always confirm: ``Was that a four-piece or forty-piece nugget?'' This adds one turn ($\sim$3 seconds) but prevents a wrong order. (4) \textbf{Cost-based confirmation threshold:} only confirm when the price difference between homophones exceeds \$5. ``Coke vs.\ Diet Coke'' (\$0 difference) $\rightarrow$ don't confirm. ``4-piece vs.\ 40-piece'' (\$9 difference) $\rightarrow$ confirm.

\textbf{Deep Dive \& Trade-offs:}
\begin{itemize}
  \item Confirmation on every item creates a tedious experience. Cost-based thresholding limits confirmations to $\sim$10\% of items.
  \item Alternative: use DTMF (keypad) input for quantities. ``Press 4 for four-piece, press 0 for forty-piece.'' Works for phone; not for drive-thru.
  \item Long-term: fine-tune STT on menu-specific audio (with emphasis on numeric terms) to reduce misrecognition rate.
\end{itemize}

\textbf{The Interviewer's Rubric:} Checking awareness of domain-specific STT failure modes and practical mitigation.

\bigskip\hrule\bigskip

\textbf{Q:} Two people are talking at the drive-thru speaker simultaneously---one ordering, one commenting to their passenger. How does the system handle overlapping speakers?

\textbf{The Ideal Answer:} Multi-speaker handling at drive-thru is a hard unsolved problem. Practical approaches: (1) \textbf{Speaker diarization:} identify the primary speaker (closest to the microphone) and filter out the secondary speaker. Modern diarization models (pyannote) can separate speakers in real-time but add 100--200ms latency---tight for the budget. (2) \textbf{Dominant speaker extraction:} without full diarization, use audio energy and direction (if using a directional microphone array) to focus on the loudest/closest speaker. Simpler, faster ($<$50ms), works 80\% of the time. (3) \textbf{Confusion detection:} if the STT output is incoherent (two interleaved conversations), the system recognizes confusion and asks: ``Sorry, I heard a couple of things. Could you repeat your order?'' (4) \textbf{Hardware solution:} directional microphone with cardioid pattern pointed at the driver-side window. Rejects speech from the passenger side by 15--20 dB. This is the most effective solution and most drive-thru installations use it.

\textbf{Deep Dive \& Trade-offs:}
\begin{itemize}
  \item Real-time speaker diarization adds 100--200ms---may break the 500ms budget if the pipeline is already tight.
  \item Directional microphones solve 90\% of the problem at hardware level, zero latency impact.
  \item For phone calls (not drive-thru), there's no directional mic. Must rely on diarization or simply accept that multi-speaker phone calls are harder.
  \item Worst case: if the system can't distinguish speakers, it processes whatever STT produces and relies on the confirmation step to catch errors.
\end{itemize}

\textbf{The Interviewer's Rubric:} Checking understanding that voice pipeline challenges are often solved at the hardware/acoustic layer, not just software.

\bigskip\hrule\bigskip

\textbf{Q:} Your LLM generates a response but the TTS pronunciation is wrong---it says ``live'' (as in ``alive'') instead of ``live'' (as in ``broadcast''). How do you control pronunciation?

\textbf{The Ideal Answer:} TTS pronunciation errors on homographs (``live,'' ``read,'' ``lead,'' ``wind'') are common. Controls: (1) \textbf{SSML (Speech Synthesis Markup Language):} wrap ambiguous words with phoneme hints: \texttt{<phoneme alphabet="ipa" ph="la\textsci v">live</phoneme>}. This requires the LLM to generate SSML-annotated text, which is feasible with prompt engineering. (2) \textbf{Pronunciation dictionary:} the TTS system maintains a custom pronunciation dictionary mapping words to phonemes in context. ``Live music'' $\rightarrow$ /la\textsci v/, ``live broadcast'' $\rightarrow$ /la\textsci v/, ``live animals'' $\rightarrow$ /la\textsci v/. Context-dependent rules. (3) \textbf{Domain-specific TTS fine-tuning:} fine-tune the TTS model on domain-specific text-audio pairs. For a drive-thru, the vocabulary is limited; pronunciation errors are rare after fine-tuning on menu items. (4) \textbf{Post-synthesis verification (expensive):} run a lightweight STT on the generated audio and compare against the input text. If mismatched, re-synthesize. Adds $\sim$200ms per turn.

\textbf{Deep Dive \& Trade-offs:}
\begin{itemize}
  \item SSML adds complexity to the LLM prompt and output parsing. Not all TTS systems fully support SSML.
  \item Pronunciation dictionaries are effective for domain-specific terms but don't generalize to novel text.
  \item TTS fine-tuning requires 5--10 hours of recorded audio in the target voice + domain text. $\sim$1 week of effort.
  \item In practice, for constrained domains (drive-thru, appointment scheduling), the vocabulary is small enough that a pronunciation dictionary covers 99\% of cases.
\end{itemize}

\textbf{The Interviewer's Rubric:} Checking awareness of TTS quality control beyond ``plug in a TTS API.''

\bigskip\hrule\bigskip

\textbf{Q:} You need to deploy this system across 5,000 drive-thru locations worldwide. How do you handle the infrastructure?

\textbf{The Ideal Answer:} Two architectural options: (1) \textbf{Cloud-centralized:} all audio streams back to regional cloud data centers. Pros: centralized GPU fleet, easy model updates, shared capacity. Cons: network latency from store to cloud (50--100ms RTT), bandwidth costs (32KB/s per stream $\times$ 5K locations), dependency on WAN. (2) \textbf{Edge deployment:} deploy a small compute node at each location (or per-cluster of locations). Run VAD + lightweight processing locally; stream only to STT/LLM in a regional data center. Pros: reduced network dependency, lower latency. Cons: 5,000 edge devices to manage, limited compute per location.

\textbf{Recommended: Hybrid.} VAD runs on a local edge device (Raspberry Pi-class, \$50 hardware). Audio streams to the nearest regional data center (8--12 regions globally) for STT$\rightarrow$LLM$\rightarrow$TTS processing. Audio transport over WebRTC with adaptive bitrate.

Regional data center sizing: 5,000 locations / 10 regions = 500 locations per region. At peak (lunch rush): 500 locations $\times$ 3 concurrent orders = 1,500 concurrent sessions per region. GPU fleet per region: $\sim$30 GPUs for STT, $\sim$5 for LLM, $\sim$15 for TTS = $\sim$50 GPUs per region, 500 GPUs total.

Network resilience: if WAN to cloud fails, the edge device plays pre-recorded prompts (``Our AI system is temporarily unavailable. A team member will take your order.'') and escalates to human.

\textbf{Deep Dive \& Trade-offs:}
\begin{itemize}
  \item Edge VAD eliminates false triggers before audio leaves the premises, saving bandwidth.
  \item Running full STT/LLM/TTS on-premise is theoretically possible (with a \$5K GPU box per location) but operationally nightmarish at 5,000 locations.
  \item Network quality varies by location. Rural locations with poor connectivity may need a larger edge compute footprint.
  \item Model updates: push new model weights to regional data centers (trivial). Push new VAD models to edge devices (non-trivial---requires OTA update infrastructure).
\end{itemize}

\textbf{The Interviewer's Rubric:} Checking infrastructure architecture thinking for geographically distributed real-time systems.

\bigskip\hrule\bigskip

\textbf{Q:} A customer has a heavy accent that drops STT accuracy to 70\% WER. How do you handle graceful degradation?

\textbf{The Ideal Answer:} Accent robustness is an equity and accessibility concern, not just a quality metric. Layers: (1) \textbf{Multi-model STT ensemble:} run two STT models in parallel (e.g., Deepgram + Whisper). Take the transcript with higher confidence. Different models have different accent strengths. (2) \textbf{Confidence-based clarification:} if STT confidence on any word drops below 60\%, ask for clarification on that specific word rather than the whole utterance: ``I got `large number three'---was the drink a Coke or a Fanta?'' (3) \textbf{Multimodal fallback:} if text-based understanding fails, route the audio to a native multimodal model (GPT-4o voice mode) that processes audio directly---it may understand accented speech that STT couldn't transcribe. (4) \textbf{Human escalation:} if confidence stays below threshold for 3 consecutive turns, offer: ``Let me connect you with a team member who can help.'' Don't let the customer struggle. (5) \textbf{Long-term:} collect and label accent-diverse training data (with speaker consent). Fine-tune STT models to improve accuracy across accent groups. Monitor WER by demographic cohort (accent, age) to ensure equitable performance.

\textbf{Deep Dive \& Trade-offs:}
\begin{itemize}
  \item Multi-model STT doubles compute cost. Only activate for low-confidence sessions.
  \item Clarification questions add friction. Limit to 1 clarification per turn.
  \item Human escalation is the safety net but defeats the purpose of automation. Track escalation rate by accent group---if one group escalates 5x more, the model needs retraining for that accent.
  \item Ethical consideration: if the system performs significantly worse for certain accent groups, it's discriminatory. Regular fairness audits on WER by accent are mandatory.
\end{itemize}

\textbf{The Interviewer's Rubric:} Checking fairness and accessibility thinking in voice AI systems.

\bigskip\hrule\bigskip

\textbf{Q:} The restaurant changes its menu seasonally. How do you update the voice agent without downtime?

\textbf{The Ideal Answer:} Menu updates propagate through three components: (1) \textbf{Dialogue context:} the system prompt and menu context (item names, prices, availability) are stored in a configuration database, not hardcoded in the model. Menu update = database record update, zero model changes. The LLM reads the current menu from context at the start of each session. (2) \textbf{STT hot-words:} the keyword boost list must be updated with new item names. Push new hot-word list to the STT service (API parameter, no model reload). (3) \textbf{TTS pronunciation dictionary:} add pronunciation entries for new menu items. Deploy as a config update to the TTS service.

Deployment: blue-green deployment at the session level. New sessions use the updated config (new menu). Existing sessions continue with the old config until they end (max 5 minutes). No mid-session menu changes to avoid confusion. Zero downtime.

Testing: before deployment, run the updated config through an automated test suite: 50 synthetic orders covering new menu items. Verify: STT recognizes new item names, LLM correctly adds them to orders, TTS pronounces them correctly, confirmation messages include correct prices.

\textbf{Deep Dive \& Trade-offs:}
\begin{itemize}
  \item Separating menu data from model weights is critical. If menu items were in the model's training data, updates would require model retraining.
  \item Hot-word lists have limits (typically 100--500 terms). Seasonal menus fit easily; restaurants with 500+ items may need tiered boosting (boost seasonal items + top-50 by popularity).
  \item The test suite should run in $<$5 minutes. Block deployment if any test fails.
  \item Rollback: if the new menu config causes issues (wrong prices, unrecognized items), revert the database record. All new sessions immediately use the old menu. Rollback time: $<$1 minute.
\end{itemize}

\textbf{The Interviewer's Rubric:} Checking operational maturity for configuration-driven updates in real-time systems.

\bigskip\hrule\bigskip

\textbf{Q:} How do you measure whether the voice agent is actually good? What metrics matter beyond latency?

\textbf{The Ideal Answer:} Voice agent quality is multi-dimensional. Metric stack: (1) \textbf{Task completion rate:} did the order get placed correctly / did the appointment get scheduled / did the issue get resolved? Ground truth: compare agent-recorded order vs.\ POS system receipt. Target: $>$90\%. (2) \textbf{Order accuracy:} for drive-thru, exact match between what the agent understood and what the customer received. Track mismatches at the item and modifier level (wrong size, missing ``no pickles''). (3) \textbf{Average handle time (AHT):} total session duration. Compare against human operator baseline. Target: within 20\% of human AHT. (4) \textbf{Escalation rate:} \% of sessions escalated to human. Target: $<$10\%. (5) \textbf{Customer satisfaction (CSAT):} post-call survey or inferred from behavioral signals (repeat usage, abandonment). (6) \textbf{Utterance understanding rate:} \% of user utterances that the system correctly interpreted (manual sampling of 500/week). (7) \textbf{Barge-in success rate:} when users interrupt, does the system correctly handle the interruption? (8) \textbf{Revenue per session:} for sales agents, compare average order value (AI agent vs.\ human). AI agents are consistent (always upsell) but may miss contextual cues.

\textbf{Deep Dive \& Trade-offs:}
\begin{itemize}
  \item Task completion rate is the ultimate metric but is expensive to measure (requires ground truth from the fulfillment system).
  \item AHT is easy to measure but can be gamed (agent that confirms excessively has low error but high AHT).
  \item CSAT surveys have low response rates ($<$5\%). Infer satisfaction from behavioral signals: repeat usage, session duration, abandonment timing.
  \item Weekly review of 500 sampled sessions by QA team provides qualitative insights that metrics miss.
\end{itemize}

\textbf{The Interviewer's Rubric:} Checking holistic quality measurement beyond engineering metrics.

\section*{End of Chapter}
\begin{itemize}
  \item \textbf{Key takeaway:} Voice-first agents are latency-engineering problems first and AI problems second. The 500ms voice-to-voice budget requires streaming at every pipeline stage---STT, LLM, and TTS cannot operate in batch mode. Barge-in handling (interruption detection and response within 200ms) is the feature that separates usable voice agents from frustrating ones. At 50K concurrent sessions, self-hosted STT and TTS are 23x cheaper than managed APIs, making infrastructure architecture as critical as model selection.
  \item \textbf{What to practice:} Draw the end-to-end latency waterfall from user speech $\rightarrow$ VAD $\rightarrow$ STT $\rightarrow$ LLM $\rightarrow$ TTS $\rightarrow$ audio playback, with millisecond budgets per stage. Identify which stage is the bottleneck and how to shave 50ms from it.
  \item \textbf{Interview drill:} Walk through a barge-in scenario: the agent is saying ``Your total comes to twelve dollars and fifty cents'' and the customer interrupts at ``twelve dollars.'' Trace every system action from VAD detection through TTS stop through STT transcription of the interruption through LLM re-generation through audio playback.
\end{itemize}
